% Generated by Listas/make_statements.py from pratica-lecture-01.md, item 19. Do not edit.
\textbf{Source.} Muñoz-Garcia \& Toro-Gonzalez (2019), \emph{Strategy and Game Theory: Practice Exercises with Answers}, Ch. 1, Exercise 16 ``Equilibrium Predictions from IDSDS Versus IDWDS'', p.~31 {[}PDF 44{]}. Reproduced verbatim (statement only; the book's worked answer follows it).



\begin{quote}

In previous exercises applying IDSDS we sometimes started finding strictly dominated strategies for the row player, while in other exercises we began identifying strictly dominated strategies for the column player (or the matrix player in games with three players). While the order of deletion does not affect the equilibrium outcome when applying IDSDS, it \emph{can} affect the set of equilibrium outcomes when we delete weakly (rather than strictly) dominated strategies.



Use the game in Fig. 1.37 to show that the order in which weakly dominated strategies are eliminated can affect equilibrium outcomes.



{\def\LTcaptype{none} % do not increment counter

\begin{longtable}[]{@{}llll@{}}

\toprule\noalign{}

Player 1 ~Player 2 & \(L\) & \(M\) & \(R\) \\

\midrule\noalign{}

\endhead

\bottomrule\noalign{}

\endlastfoot

\(U\) & 2, 1 & 1, 1 & 0, 0 \\

\(C\) & 1, 2 & 3, 1 & 2, 1 \\

\(D\) & 2, -2 & 1, -1 & -1, -1 \\

\end{longtable}

}



\emph{Fig. 1.37} Normal-form game

\end{quote}

